\section{Real incidence estimates}
\label{ci:section}

We give the incidence arguments used in the proof.  All constants in this
section are absolute.  The elementary algebraic facts used below are the
following forms of B\'ezout's Theorem: two relatively prime polynomials of
degrees $a,b$ in three variables have at most $ab$ distinct common line
components, and a polynomial of degree $a$ has at most $a$ zeros on a line
unless it vanishes identically on that line.  We also use unique
factorization, the implicit function theorem, and the Borsuk--Ulam Theorem
in its standard form that every continuous odd map $S^q\longrightarrow
\mathbb R^q$ has a zero \cite[Corollary~2B.7]{hatcher2002}.
In the ruled-surface argument we use the basic
projective facts that a projective morphism has closed image, that
varieties have finitely many irreducible components, and that a proper
closed subset of an irreducible curve is finite.  We use the dimension
and generic-fiber-dimension theorems, and extension of derivations across
separable algebraic field extensions.  The incidence theorems and the
special assertions about ruled surfaces are proved below.
The historical sources for the two principal incidence bounds
are \cite{szemereditrotter1983,guthkatz2015}.

\subsection{The planar incidence bound}

We use the crossing inequality of Ajtai, Chv\'atal, Newborn, and
Szemer\'edi and of Leighton
\cite{ajtai1982crossing,leighton1983}, following Sz\'ekely's
crossing-number proof of the incidence estimate~\cite{szekely1997}.

\begin{lemma}[Crossings]\label{ci:crossing}
Suppose a simple graph with $v\ge1$ vertices and $e\ge4v$ edges is drawn in
the plane, no edge goes through a vertex other than its endpoints, and
two edges with a common endpoint have no other intersection.  If $X$
counts intersecting pairs of edges away from their endpoints, then
$X\ge e^3/(64v^2)$.
\end{lemma}

\begin{proof}
Deleting at most one edge for each intersecting pair produces a planar
simple graph.  Euler's formula therefore gives $e\le3v+X$.  Select each
vertex independently with probability $p$ and keep the induced drawing.
An edge survives with probability $p^2$ and an intersecting pair with
probability $p^4$, because its four endpoints are different.  Taking
expectations in the same inequality gives
$p^2e\le3pv+p^4X$.  Set $p=4v/e$ and rearrange.
\end{proof}

\begin{theorem}[Szemer\'edi--Trotter]\label{ci:st}
For finite sets $Q$ of $n$ points and $\mathcal L$ of $L$ distinct real
lines in the plane,
\[
 I(Q,\mathcal L)\le C\bigl(n^{2/3}L^{2/3}+n+L\bigr).
\]
Consequently the number of points incident to at least $k\ge2$ lines is
at most $C(L^2/k^3+L/k)$, and the number of lines containing at least
$k\ge2$ points of an $n$-point set is at most
$C(n^2/k^3+n/k)$.
The same assertions apply to finite point--line configurations in
$\mathbb R^s$ and in real projective space.
\end{theorem}

\begin{proof}
The assertion is immediate if $n=0$ or $L=0$, so assume both are positive.
Join consecutive points of $Q$ on each line of $\mathcal L$.  The
resulting graph is simple: a pair of distinct points determines just
one line.  If $e$ is its number of edges, then $e\ge I-L$.  Edges on
the same line have disjoint interiors.  Edges on two different lines
cross at most once, so their crossing count is at most $\binom L2$.
Lemma~\ref{ci:crossing} gives either $e<4n$ or
$e^3\le32n^2L^2$, proving the incidence bound.  Concurrent crossings
cause no problem: count pairs of intersecting edges, or separate the
crossings in a sufficiently small disk without changing the bound.

For a set of $s$ points each incident to at least $k$ lines, the bound
gives
$ks\le C(s^{2/3}L^{2/3}+s+L)$.  If $k$ exceeds a fixed constant,
absorb the term $Cs$ and compare the two remaining terms, obtaining
$s\le C'(L^2/k^3+L/k)$.  For bounded $k\ge2$, count pairs of lines:
$s\le\binom L2$, which gives the same assertion after increasing the
constant.  Interchanging points and lines by planar projective duality
proves the rich-line statement.  One can instead apply the incidence
bound directly to the finite collection of rich lines, using that two
points determine at most one line for bounded $k$.

For a finite configuration in higher-dimensional affine space, choose
a linear projection to $\mathbb R^2$ which is injective on its points
and sends its lines to distinct lines.  Each forbidden coincidence is
a proper algebraic condition on the projection, and finitely many such
conditions can be avoided.  Existing incidences survive, which is all
that an upper bound requires.  For projective data first choose an
affine chart containing every specified point and in which no specified
line is contained in the hyperplane at infinity.  Apply the affine
argument.
\end{proof}

\subsection{Polynomial bisection and flat lines}

Polynomial bisection is the Stone--Tukey ham-sandwich construction
\cite{stonetukey1942}, as used in Guth--Katz polynomial partitioning
\cite{guthkatz2015}.

\begin{lemma}[Polynomial bisection]\label{ci:partition}
Given a finite set $S\subset\mathbb R^3$ and an integer $r\ge1$, there
is a nonzero real polynomial $f$ of degree at most $C r$ such that
$\mathbb R^3\setminus Z(f)$ is partitioned into at most $C r^3$ open
sets, each containing at most $C|S|/r^3$ points of $S$.  Every line not
contained in $Z(f)$ meets at most $\deg f+1$ of these sets.  The sets
need not be connected.
\end{lemma}

\begin{proof}
For $m<\binom{D+3}{3}$ finite-volume bounded measurable sets in
$\mathbb R^3$, normalize the coefficients of a degree-at-most-$D$
polynomial to lie on their unit sphere.  Its signed-volume imbalances
on the $m$ sets form a continuous odd map to $\mathbb R^m$.  Continuity
follows by dominated convergence, since the zero set of a nonzero
polynomial has Lebesgue measure zero.  The latter fact follows by
induction on the number of variables and Fubini's Theorem.  Restrict
the coefficient sphere to an $(m+1)$-dimensional linear subspace and
apply Borsuk--Ulam.  We obtain a polynomial bisecting all $m$ sets.

Replace each point in each of $m$ finite sets by a disjoint small ball
of the same radius within that set, apply this assertion, and let the
radius tend to zero.  A convergent subsequence of normalized coefficient
vectors gives a nonzero polynomial for which each strict sign contains
at most half the points of each finite set.  Indeed, a strict positive
value at a point persists on its small ball and for nearby coefficient
vectors; more than half of the points of either sign would contradict
the volume bisection.

Iterate this finite bisection $J$ times, bisecting all the at most
$2^{j-1}$ sign classes at stage $j$.  A degree $C2^{j/3}$ suffices at
that stage.  The product $f$ of these $J$ polynomials has degree at
most $C2^{J/3}$, and every strict sign class has at most $2^{-J}|S|$
points.  Choose $2^J$ comparable to $r^3$.  A line outside $Z(f)$ is
cut into at most $\deg f+1$ intervals by its intersections with
$Z(f)$; all the factors have constant sign on each interval.  This
proves the last assertion as well.
\end{proof}

The critical-line and flat-line tests below are part of the algebraic
incidence method of Guth--Katz and Elekes--Kaplan--Sharir
\cite{guthkatz2010kakeya,elekeskaplansharir2011}.
For a square-free real polynomial $f$ of degree $D$, call a point of
$Z(f)$ \emph{critical} if $\nabla f=0$.  A noncritical point is
\emph{flat} if its second fundamental form vanishes.  A line in $Z(f)$
is critical if all its points are critical, and is flat if it is not
critical and all its noncritical points are flat.

\begin{lemma}[Critical and flat lines]\label{ci:flat}
For a square-free polynomial $f$ of degree $D$:
\begin{enumerate}
\item A point lying on three distinct lines contained in $Z(f)$ is
critical or flat.
\item There are at most $D(D-1)$ critical lines.
\item A line containing more than $D$ critical points is critical.
A line containing more than $3D$ flat points is flat.
\item At most $3D^2$ flat lines fail to lie in a plane component of
$Z(f)$.
\end{enumerate}
\end{lemma}

\begin{proof}
At a noncritical point all contained lines are tangent to the smooth
surface.  The second fundamental form vanishes in their directions.
A homogeneous quadratic on a two-dimensional vector space which
vanishes in three distinct projective directions vanishes identically.
This proves the first assertion.

Square-freeness in characteristic zero implies that a generic constant
linear combination of the first derivatives of $f$ shares no
irreducible factor with $f$.  All critical lines are common line
components of their zero sets, so B\'ezout gives the second assertion.
Restriction of $f$ and its derivatives to a line proves the assertion
about many critical points.

Here is an explicit polynomial test for flatness.  With $H_f$ denoting
the Hessian matrix and $e_1,e_2,e_3$ the coordinate vectors, put
\[
 T_j=\bigl(H_f(e_j\mathbin\times\nabla f)\bigr)
                  \mathbin\times\nabla f,\qquad j=1,2,3.
\]
The nine scalar components have degree at most $3D$.  At a regular
point the vectors $e_j\times\nabla f$ span the tangent plane.  Thus
all $T_j$ vanish exactly when $H_f$ maps every tangent vector into
the normal line, which is exactly vanishing of the second fundamental
form.  If a line contains more than $3D$ flat points, $f$ and all these
components vanish identically on it.  It has a regular point, so it is
a flat line.

It remains to bound flat lines on nonplanar components.  Fix an
irreducible real factor $g$ of $f$ which contains a flat line and is
not linear.  That line has a point regular for $f$, and consequently
$g$ has a smooth real point.  If every component of every $T_j$ were
divisible by $g$, the second fundamental form would vanish on a
neighborhood of that real point in $Z(g)$.  The differential of the
unit normal would vanish there.  The normal would be constant on a
smaller connected neighborhood, and that neighborhood would be
contained in its tangent plane.  The implicit function theorem then
makes it an open subset of that plane.  Restricting $g$ to the plane
shows that its plane equation divides $g$, contrary to irreducibility
and $\deg g>1$.  Therefore some component $T$ is not divisible by
$g$, and B\'ezout bounds their common lines by $3D\deg g$.

Every flat line belongs to exactly one irreducible factor of $f$:
membership in two factors would make the entire line critical.
Summing the preceding bound over the nonplanar factors gives at most
$3D^2$ flat lines.  Factors having no smooth real point contain no
flat line and do not need a separate argument.
\end{proof}

\subsection{The higher-richness line theorem}

\begin{theorem}[Guth--Katz, higher richness]\label{ci:rich-lines}
Let $\mathcal L$ be $L$ distinct lines in $\mathbb R^3$, with at most
$B_0$ in any affine plane.  For every integer $k\ge3$, the number
$R_k$ of points lying on at least $k$ lines satisfies
\[
 R_k\le C\left(\frac{L^{3/2}}{k^2}
             +\frac{LB_0}{k^3}+\frac Lk\right).
\]
There is no assumption about lines on quadrics in this theorem.
\end{theorem}

\begin{proof}
We first prove a uniform version.  Let $S$ be a set of $s$ points each
incident to between $k$ and $2k$ lines.  Suppose at least $L/100$ lines
each contain at least $sk/(100L)$ points of $S$.  We will prove the
stated bound for $s$.

Fix a sufficiently small absolute $\varepsilon>0$.  We may assume
\begin{equation}\label{ci:large-s}
 s>A\left(L^{3/2}k^{-2}+Lk^{-1}\right),
\end{equation}
where $A$ will be chosen sufficiently large.  By
Theorem~\ref{ci:st}, this assumption implies $s\le C L^2k^{-3}$.
Choose
\[
 r=\left\lfloor\theta L^2s^{-1}k^{-3}\right\rfloor,
\]
where the constant $\theta$ is sufficiently large.  Thus $r\ge1$
and $r$ is comparable to the expression inside the floor.
Lemma~\ref{ci:partition} gives a polynomial $f$ of degree
$D\le Cr$.  We claim that at most $\varepsilon s$ points of $S$ lie
outside $Z(f)$.

Otherwise, among its $O(r^3)$ sign cells, at least $c_\varepsilon r^3$
contain at least $c_\varepsilon s/r^3$ points of $S$.  To verify this,
discard cells with fewer than a small fixed multiple of
$\varepsilon s/r^3$ points and use the uniform upper bound
$Cs/r^3$ for the others.  Each line meets at most $D+1=O(r)$ cells.
One of the retained cells is therefore met by at most
$C_\varepsilon L/r^2$ lines.  Apply Theorem~\ref{ci:st} inside it:
\[
 \frac{s}{r^3}
 \le C_\varepsilon\left(\frac{L^2}{r^4k^3}
                         +\frac L{r^2k}\right).
\]
Substitution of $r$ gives
\[
 s\le C_\varepsilon\left(\theta^{-1}s
                   +\theta L^3s^{-1}k^{-4}\right).
\]
Choose $\theta$ first so that the first term can be absorbed.  The
resulting inequality $s\le C_\varepsilon\sqrt\theta L^{3/2}k^{-2}$
contradicts \eqref{ci:large-s} when $A$ is sufficiently large.  This
proves the claim.

Replacing $f$ by its square-free part does not change its zero set
and can only lower its degree.  The same estimates and
\eqref{ci:large-s} give
\begin{equation}\label{ci:small-degree}
 D\le C\theta L^2s^{-1}k^{-3},\qquad
 \frac D{sk/L}\le\frac{C\theta}{A^2},\qquad
 \frac D{\sqrt L}\le\frac{C\theta}{Ak}.
\end{equation}
Write $\mathcal L_Z$ for the lines contained in $Z(f)$, and $S_*$ for
the points of $S$ lying on at least three lines of $\mathcal L_Z$.
A point of $(S\cap Z(f))\setminus S_*$ lies on at least $k-2$ lines
outside $\mathcal L_Z$.  Each such line has at most $D$ points on
$Z(f)$.  Hence
\[
 |S\setminus S_*|\le\varepsilon s+\frac{LD}{k-2}
 \le\left(\varepsilon+\frac{3C\theta}{A^2}\right)s.
\]
Choose $A$ so large that this is at most $2\varepsilon s$.

Let $\mathcal L_0$ be the at least $L/100$ lines in the uniformity
assumption.  If such a line contains fewer than $sk/(200L)$ points
of $S_*$, it contains at least $sk/(200L)$ points of $S\setminus S_*$.
But this latter set has at most $2k|S\setminus S_*|\le4\varepsilon sk$
incidences.  Thus at most $800\varepsilon L$ lines of
$\mathcal L_0$ fail to contain $sk/(200L)$ points of $S_*$.  With
$\varepsilon$ small enough, at least $L/200$ lines succeed.

All points of $S_*$ are critical or flat by Lemma~\ref{ci:flat}.
On each successful line one of these two kinds occurs at least
$sk/(400L)$ times.  By \eqref{ci:small-degree} this is more than
$3D$, on enlarging $A$.  Every successful line is consequently a
critical line or a flat line.  Lemma~\ref{ci:flat} and the last bound
in \eqref{ci:small-degree} show that fewer than $L/400$ of them are
critical or flat on nonplanar components.  At least $L/400$ therefore
lie in plane components of $Z(f)$.  There are at most $D$ such
components, so one contains at least $L/(400D)$ lines.  It follows
that
\[
 B_0\ge\frac L{400D}\ge c\frac{sk^3}{L},
 \qquad s\le C LB_0k^{-3}.
\]
Together with the alternative to \eqref{ci:large-s}, this proves the
uniform version with an absolute constant $C_0$.

We remove the uniformity assumption by induction on $L$, keeping this
same constant $C_0$.  Let $\mathcal L_1$ consist of the lines containing
at least $sk/(100L)$ points of $S$.  If $|\mathcal L_1|\ge L/100$,
we have already finished.  Otherwise the other lines supply fewer
than $sk/100$ incidences, and at least $9s/10$ points of $S$ lie on
at least $9k/10$ lines of $\mathcal L_1$.  Partition these points
according as their multiplicity in $\mathcal L_1$ is at least $k$ or
less than $k$.  One class has at least $9s/20$ points.  Its
multiplicities lie between $k_1$ and $2k_1$, where either $k_1=k$ or
$k_1=\lceil9k/10\rceil\ge3$.  The induction hypothesis applies with
$L_1=|\mathcal L_1|<L/100$ and the unchanged plane cap $B_0$.
Consequently
\[
 s\le\frac{20}{9}C_0
 \left(\frac{L_1^{3/2}}{k_1^2}
                  +\frac{L_1B_0}{k_1^3}+\frac{L_1}{k_1}\right)
 \le C_0\left(\frac{L^{3/2}}{k^2}
                  +\frac{LB_0}{k^3}+\frac Lk\right),
\]
because $(20/9)(1/100)(10/9)^3<1$.  Empty configurations start the
induction.  Finally divide all $k$-rich points into multiplicity
intervals $[2^jk,2^{j+1}k)$ and sum this uniform-multiplicity bound.
The three resulting geometric series converge.
\end{proof}

\begin{corollary}[Planar motions with many matches]\label{ci:planar-motions}
Let $U,V\subset\mathbb R^2$ have respective cardinalities $u,v$.
For $3\le k\le\min(u,v)$, at most $C(uv)^{3/2}/k^2$ Euclidean
isometries $g$ have $|g(U)\cap V|\ge k$.
\end{corollary}

\begin{proof}
There are finitely many isometries with at least two matches: two
distinct ordered matches determine at most two isometries.  Separate
the two orientation classes and reflect one cloud for the reversing
class.  A generic rotation of $V$ makes $I+R$ invertible for the
linear part $R$ of every counted direct isometry.  Write
$J(x_1,x_2)=(-x_2,x_1)$.  In its planar Cayley chart a direct
isometry is represented by a point $(t,c)\in\mathbb R\times\mathbb R^2$,
and a match $x\mapsto y$ is represented by the line
\begin{equation}\label{ci:cayley-line}
 c=\frac{x-y}{2}-tJ\frac{x+y}{2}.
\end{equation}
Indeed, putting $Z=tJ$, the match equation is equivalent to
$(I+Z)y=(I-Z)x-2c$, and $(I+Z)^{-1}(I-Z)$ is orthogonal of
determinant one.  Conversely every direct isometry with $I+R$
invertible has this representation.

The $uv$ lines are distinct: their intercept and slope determine
$x-y$ and $x+y$, hence both endpoints.  A plane in $(t,c)$-space
has an equation $a\cdot c+bt=d$.  If $a=0$ it contains no such
line.  Otherwise, for a line \eqref{ci:cayley-line} contained in
the plane,
\[
 a\cdot(x-y)=2d,\qquad a\cdot J(x+y)=2b.
\]
With $x$ fixed these determine $y$ uniquely, because $a$ and
$J^Ta$ are independent; the same holds with $y$ fixed.  Thus the
plane cap is $B_0\le\min(u,v)$.  Apply
Theorem~\ref{ci:rich-lines}.  Its middle term is bounded by its
first because $B_0\le\sqrt{uv}\le k\sqrt{uv}$, and its last is
bounded by its first because $k\le\min(u,v)\le\sqrt{uv}$.
The bounded orientation split completes the proof.
\end{proof}

\subsection{A ruled-surface criterion with a degree bound}

This is the flecnode approach to ruledness used in
\cite[Section~3, Proposition~3.2 and Corollary~3.3]{guthkatz2015}.
The direct resultant calculation below gives degree at most $17d-24$;
the sharper flecnode degree used in that source is not needed here.

All varieties in the next two lemmas are over $\mathbb C$, unless a
real structure is explicitly mentioned.  A projective surface is
\emph{ruled} if every point belongs to a projective line contained in
the surface.  The statement is geometric: a real sphere is ruled
over $\mathbb C$ even though it contains no real line.  A
\emph{regulus} in the real incidence statements is the affine part
of a smooth real projective quadric with real lines.

\begin{lemma}[A polynomial detecting ruledness]\label{ci:flecnode}
An irreducible surface of degree $d$ in $\mathbb C^3$ which is not
geometrically ruled contains at most $17d^2$ affine lines.
\end{lemma}

\begin{proof}
Planes are ruled.  An irreducible complex projective quadric has
rank three or four: rank three gives a cone over a conic, and in
rank four a linear coordinate change gives $X_0X_3=X_1X_2$,
whose two rulings follow by fixing the ratio in either pair of
factors.  Hence we may assume $d\ge3$.
Write its irreducible defining polynomial as $F$.  After a linear
coordinate change, $A=F_z$ is not identically zero.  Put $B=F_x$,
$C=F_y$ and form the tangent vector
\[
 V(u,v)=(Au,Av,-Bu-Cv).
\]
The binary forms
\[
 q(u,v)=D^2F[V,V],\qquad c(u,v)=D^3F[V,V,V]
\]
are obtained by differentiating $F$ only, with the displayed vector
held fixed.  Their coefficient degrees are at most $3d-4$ and
$4d-6$.  Let $R$ be their homogeneous binary resultant.  The
Sylvester determinant is homogeneous of degree three in the
quadratic coefficients and degree two in the cubic coefficients;
therefore
\begin{equation}\label{ci:resultant-degree}
 \deg R\le3(3d-4)+2(4d-6)=17d-24.
\end{equation}
The resultant vanishes precisely when the two binary forms have a
common projective root; this includes identically zero forms.

Every affine line in $Z(F)$ lies in $Z(R)$.  At a point where
$A\ne0$, use its tangent direction, for which all directional
derivatives vanish.  At a point where $A=0$, the two forms become
\[
 q=F_{zz}(Bu+Cv)^2,\qquad
 c=-F_{zzz}(Bu+Cv)^3,
\]
so they again have a common projective root.  This also handles a
line lying entirely in $Z(F_z)$.

We next prove that $F\mid R$ implies ruledness.  Let
$K=\mathbb C(Z(F))$.  This field is a finite separable extension of
$\mathbb C(x,y)$.  Extend $\partial_x,\partial_y$ to an algebraic
closure of $K$, where they commute, and regard $z$ as an algebraic
function of $x,y$.  Write
\[
 r=z_{xx},\qquad s=z_{xy},\qquad t=z_{yy}.
\]
If this Hessian is zero, $D=\partial_x$ satisfies
$D^2x=D^2y=D^2z=0$, and the formal argument below already gives a
line.  Otherwise a constant invertible linear change of $x,y$
ensures $t\ne0$.  Thus the binary quadratic of second derivatives
does not vanish in the infinite-slope direction.

Since $F\mid R$, there is a finite common tangent slope $a$ in the
algebraic closure of $K$.  Put $w=z_x+az_y$.  Implicit
differentiation, keeping $a$ fixed in the directional derivatives,
gives
\begin{align*}
 D^2F[(1,a,w),(1,a,w)]&=-F_z Q(a),\\
 D^3F[(1,a,w),(1,a,w),(1,a,w)]
 &=-F_z P(a)-3Q(a)(F_{xz}+aF_{yz}+wF_{zz}),
 \end{align*}
where
\[
 Q(a)=r+2sa+ta^2,\qquad
 P(a)=z_{xxx}+3az_{xxy}+3a^2z_{xyy}+a^3z_{yyy}.
\]
Consequently $Q(a)=P(a)=0$.

Set $D=\partial_x+a\partial_y$.  If $a$ is a simple root of $Q$,
differentiating $Q(a)=0$ gives
\[
 0=D(Q(a))=P(a)+2(s+ta)Da,
\]
so $Da=0$.  If it is a repeated root, then $s=-ta$ and $r=ta^2$.
Use the mixed-partial identities $r_y=s_x$, $s_y=t_x$:
\[
 t_x=-t_ya-ta_y,\qquad
 t_ya^2+2taa_y=-t_xa-ta_x
             =t_ya^2+taa_y-ta_x.
\]
Their difference gives $t(a_x+aa_y)=0$, so $Da=0$ in this case
as well.  It follows that
\[
 Dx=1,\quad Dy=a,\quad Dz=w,\qquad
 D^2x=D^2y=D^2z=0.
\]
The formal Taylor map
$\exp(TD)=\sum_{j\ge0}T^jD^j/j!$ is an algebra homomorphism in
characteristic zero: the product rule proves this coefficient by
coefficient.  Apply it to $F(x,y,z)=0$ to obtain
\[
 F(x+T,y+aT,z+wT)=0.
\]
Thus the geometric generic point of the surface lies on a contained
line.  The set of lines on its projective closure is a closed subset
of the Grassmannian: substituting the line parameterization in the
homogenized equation gives polynomial equations for it.  The
projective point--line incidence variety has closed image in the
surface.  Since its image contains the geometric generic point,
it is the whole surface.  This proves ruledness.

If the surface is not ruled, $F$ and $R$ are therefore relatively
prime, and B\'ezout with \eqref{ci:resultant-degree} bounds their
common lines by $17d^2$.
\end{proof}

\subsection{Intersections among generators}

\begin{lemma}[Ruled-surface structure]\label{ci:ruled-structure}
Let $S$ be an irreducible ruled projective surface over $\mathbb C$
which is neither a plane nor a smooth quadric.  Its lines have a
unique one-dimensional irreducible family $\Gamma$, called the
generator family, whose members cover $S$.  A generic point of $S$
belongs to exactly one generator.  There is at most one point lying
on infinitely many contained lines, and at most two contained lines
meeting infinitely many other contained lines away from such points.
Every line outside $\Gamma$ is one of these at most two lines.
\end{lemma}

\begin{proof}
The algebraic set of lines on a nonplanar irreducible surface has
dimension at most one.  Otherwise its incidence variety would have
dimension at least three, and through a generic smooth point of
$S$ there would be infinitely many contained lines.  All would lie
in the tangent plane.  Infinitely many distinct lines through that
point fill the tangent plane algebraically, making this plane a
component of $S$, a contradiction.  Each curve component of the
line variety covers $S$: its union is closed by projectivity, and
infinitely many distinct lines cannot be contained in a proper
closed subset of an irreducible surface.  Indeed that subset has
finitely many curve components.  All other components of the line
variety are isolated points and thus represent finitely many lines.

Suppose a generic point had at least two generators. Fix a generic
generator $\ell$ from a component $\Gamma_1$. A dense set of points
of $\ell$ then belongs to a second generator distinct from $\ell$.
These second generators are infinitely many, since each meets
$\ell$ in only one point. Some curve component $\Gamma_j$ therefore
has infinitely many members meeting $\ell$, and closedness makes
every member meet $\ell$. For each $j$, the set of members of
$\Gamma_1$ meeting every member of $\Gamma_j$ is closed: it is the
intersection of the closed meeting loci for the individual lines.
These finitely many closed sets cover a dense subset of
$\Gamma_1$. Irreducibility makes one of them all of $\Gamma_1$.
Call the corresponding component $\Gamma_2$. This includes the
case $\Gamma_2=\Gamma_1$; the diagonal alone was not counted as a
second generator.

If three generic members of $\Gamma_1$ are pairwise skew, they
determine a unique smooth quadric, and their common transversals
are exactly its opposite ruling.  To see this directly, take the
first two lines as complementary coordinate two-planes in
$\mathbb C^4$ and the third as the graph of an invertible linear
map between them.  A transversal is the projectivization of the
span of $u$ and its image, with $[u]\in\mathbb P^1$; their union
is a nonsingular quadric.  Hence $\Gamma_2$ would cover that quadric
and $S$ would equal it.

If generic pairs in $\Gamma_1$ are not skew, every pair meets,
because the meeting condition is closed in
$\Gamma_1\times\Gamma_1$.  A family of pairwise meeting lines is
coplanar or concurrent.  Indeed, fix two meeting at $p$ and spanning
a plane $H$.  Every line meeting both is contained in $H$ or passes
through $p$; irreducibility of the family forces one alternative
throughout.  The first makes $S$ a plane.  The second makes it a
cone with vertex $p$.  Such a nonplanar cone contains no line
$\ell$ avoiding $p$: all the joins from $p$ to the points of
$\ell$ would give the entire plane spanned by $p$ and $\ell$ in
$S$.  It consequently has a unique generator through a generic
point, contradicting the assumption of a second generator.
This proves generic uniqueness.  It also proves uniqueness of
$\Gamma$, since every curve component covers $S$.

Call a point or line \emph{exceptional} if it has the corresponding
infinite-incidence property in the statement.  If $p$ is an
exceptional point, the closed subfamily of lines through it contains
the whole generator curve.  Thus all generators pass through $p$.
There cannot be two such points: a generator through both would
be their unique joining line.  Hence there is at most one
exceptional point.

An exceptional line meets every generator, since infinitely many
generators meet it and $\Gamma$ is irreducible.  Three pairwise skew
exceptional lines would force $S$ into their quadric as above.
Two exceptional lines cannot meet: all generators meeting both
would lie in their plane or pass through their common point.
The former makes $S$ planar, while the latter makes the common
point exceptional and prevents either line from meeting infinitely
many generators away from it.  Thus there are at most two
exceptional lines.  Finally, a contained line outside $\Gamma$
has all its points covered by generators.  Infinitely many
generators meet it at distinct points, and removing the possible
one exceptional point does not change this fact.  It is therefore
exceptional.
\end{proof}

\begin{lemma}[Counting on ruled surfaces]\label{ci:ruled-count}
Let a union of real absolutely irreducible geometrically ruled
surfaces have total degree $d\ge1$, and let $\mathcal L$ consist of
$M\ge1$ distinct real affine lines contained in that union.  If at
most $s$ lines lie in any plane or regulus, then the number of
points on at least two lines of $\mathcal L$ is at most
$5dM+sM/2$.
\end{lemma}

\begin{proof}
First suppose the projective closure of every irreducible component
is neither a plane nor a smooth projective quadric.  Across all
the components there are at most $d$
exceptional points and at most $2d$ exceptional lines by
Lemma~\ref{ci:ruled-structure}.  Intersections involving an
exceptional line cost at most $2dM$ points.  The exceptional points
themselves cost at most $d$.

Fix any remaining line $\ell$.  Choose a generic complex plane
$H$ through it and restrict the square-free equation of the entire
union to $H$.  Divide out the entire power of the linear factor
defining $\ell$.  The residual plane curve $C_H$ has degree at
most $d-1$ and does not contain $\ell$.  Dividing out the entire
power is necessary when $\ell$ is a singular line of a component.

We claim that every remaining intersection $q=\ell\cap\ell'$
belongs to $C_H$.  The nonexceptional line $\ell'$ belongs to the
generator curve of some component $S_j$.  Take $H$ not containing
$\ell'$.  If only finitely many generators of $S_j$ meet $\ell$,
then the map sending a generator to its intersection with $H$
is regular near $\ell'$.  On a dense open subset its values lie
in $C_H$.  As $C_H$ is closed, its value at $\ell'$ is also in
$C_H$, which says $q\in C_H$.  This argument applies even at a
singular point of the generator curve.

The alternative of infinitely many generators meeting $\ell$
cannot occur at the present $q$.  If $\ell\subset S_j$, either
these meetings range through infinitely many nonexceptional
points, making $\ell$ exceptional, or all generators pass through
a single exceptional point on $\ell$.  In the second case every
generator other than $\ell$ meets it only at that point.  If
$\ell\not\subset S_j$, the finite set $\ell\cap S_j$ contains
all the meetings; irreducibility then forces all generators
through one point of this finite set, which is exceptional.
Both cases contradict the choice of $q$.  This proves the claim.
B\'ezout in $H$ now bounds the number of such $q$ on $\ell$ by
$d-1$.  Summing and restoring the exceptions bounds the total by
$(d-1)M+2dM+d\le4dM$.

For the general union, extract all lines in each plane component
or component with smooth projective quadric closure, one component
at a time.  A smooth real projective quadric having real lines
gives a regulus; a quadric without real
lines contributes nothing.  If an extracted group has $a_j$
lines, then $a_j\le s$ and its internal intersections number at
most $a_j^2/2\le s a_j/2$.  Every line left after its extraction
is not contained in that component, so it meets it in at most
the degree of the component points.  All such cross intersections
cost at most $dM$.  The remaining components and lines are covered
by the preceding $4dM$ estimate.  Since the extracted groups are
disjoint, their internal contributions sum to at most $sM/2$.
\end{proof}

\subsection{The two-rich line theorem}

The sampling and degree-reduction argument follows the architecture of
the two-rich-point proof in \cite[Section~3]{guthkatz2015}.

\begin{theorem}[Two-rich points]\label{ci:two-rich}
For $m$ distinct real affine lines, with at most $s\ge1$ in any
plane or regulus, the number $P_2$ of points lying on at least two
lines satisfies
\[
 P_2\le K\bigl(m^{3/2}+ms\bigr)
\]
for an absolute constant $K$.
\end{theorem}

\begin{proof}
Write $H(m,s)=m^{3/2}+ms$ and induct on $m$, keeping the same
parameter $s$ throughout.  Choose $K$ larger than all the absolute
constants required below.  The case $m\le1$ is immediate.
Suppose, towards a contradiction, that $P_2>KH(m,s)$.  Repeatedly
delete a line having fewer than
\[
 T=\frac K4(\sqrt m+s)
\]
distinct intersection points with other currently retained lines.
Each deletion loses fewer than $T$ two-rich points.  The remaining
set of $M\le m$ lines therefore has more than $3KH(m,s)/4$
two-rich points and is nonempty.  Every remaining line has at
least $T$ distinct intersections with the other remaining lines.

Select each remaining line independently with probability $64/K$.
For each fixed remaining line choose one witness line at every
distinct intersection point.  These witnesses are distinct.  The
number $X$ of selected witnesses is a binomial random variable
with mean at least $16\sqrt m$.  The elementary binomial tail
bound $\Pr(X<\mathbb EX/2)\le\exp(-\mathbb EX/8)$ gives
\[
 \Pr(X<8\sqrt m)\le e^{-2\sqrt m}.
\]
For completeness, this tail bound follows by applying Markov's
inequality to $e^{-\lambda X}$ and the inequality
$1+p(e^{-\lambda}-1)\le\exp(p(e^{-\lambda}-1))$, then taking
$\lambda=\log2$; the resulting exponent is at most
$-\mathbb EX/8$.  A union bound over all remaining lines costs at
most $me^{-2\sqrt m}\le e^{-2}<1/4$.  Markov's inequality also
shows that with probability at least $1/2$ the number $q$ of
selected lines is at most $128M/K$.  There is therefore a choice
with both properties.  It necessarily has $q\ge1$.

Set $D=\lceil\sqrt{6q}\rceil$.  A polynomial of degree at most
$D$ has $\binom{D+3}{3}$ coefficients, and vanishing on a line
imposes at most $D+1$ linear conditions.  Since
$(D+3)(D+2)>6q$, there is a nonzero real polynomial vanishing on
every selected line.  Moreover
\[
 D<4\sqrt q,\qquad D^2\le2048m/K,
\]
and $D<8\sqrt m$ for our choice of $K$.  Each remaining line
contains at least $8\sqrt m$ distinct points on the selected
lines, so the polynomial vanishes identically on all $M$ remaining
lines.  Replace it by its square-free part.

Factor the polynomial over $\mathbb C$.  Retain the real,
absolutely irreducible, geometrically ruled factors.  At most
$b\le18D^2$ of the $M$ lines fail to lie on their union.  Indeed,
Lemma~\ref{ci:flecnode} bounds the lines on each nonruled factor
of degree $d_i$ by $17d_i^2$.  A real line on a factor which is
not real up to a scalar also lies on its distinct conjugate;
B\'ezout bounds such lines by $d_i^2$.  Summing these bounds and
using $\sum d_i\le D$ proves the assertion.  Thus
\[
 b\le36864m/K\le m/4.
\]

Every one of these $b$ lines meets the retained ruled union in at
most $D$ points, since it is not contained in that union.
Lemma~\ref{ci:ruled-count} bounds intersections among the ruled
lines by $5Dm+sm/2$.  The induction hypothesis, with the original
cap $s$, bounds intersections among the $b$ other lines by
$KH(b,s)$.  Altogether the remaining configuration has at most
\[
 6Dm+sm/2+KH(b,s)
 \le(49+K/4)H(m,s)
\]
two-rich points.  Here $D<8\sqrt m$ and
$H(b,s)\le H(m,s)/4$ were used.  Taking $K>98$, in addition to
the preceding requirements, contradicts the lower bound
$3KH(m,s)/4$.  This completes the induction.
\end{proof}

\subsection{The planar distinct-distance estimate}

\begin{lemma}[Concentration of planar Cayley lines]\label{ci:cayley-caps}
For two finite planar sets $U,V$ of sizes $u,v$, the $uv$ lines
\eqref{ci:cayley-line} have at most $\min(u,v)$ members in a plane
and at most $6u+2v$ members in a regulus.
\end{lemma}

\begin{proof}
Distinctness and the plane bound were proved in
Corollary~\ref{ci:planar-motions}.  Fix $x\in U$ and allow the
other endpoint $y$ to range over all of $\mathbb R^2$.  These lines
partition $(t,c)$-space: at every point their equation uniquely
determines
\[
 y=(I+tJ)^{-1}\bigl((I-tJ)x-2c\bigr).
\]
Their nonparallel directions and lack of intersections mean that
they are pairwise skew even projectively.  The polynomial vector
field
\[
 W_x(t,c)=\bigl(1+t^2,(tI+J)(c-x)\bigr)
\]
is nowhere zero over $\mathbb R$ and tangent to this family.
Indeed along its line through $(t,c)$ the derivative of $c$ with
respect to $t$ is $(tI+J)(c-x)/(1+t^2)=-J(x+y)/2$.

Let a regulus have irreducible quadratic equation $f=0$.  Every
line in this fixed-$x$ family contained in the regulus is a common
line of $f$ and $W_x\cdot\nabla f$.  The latter has degree at
most three.  If there are at least seven such lines, B\'ezout
forces
\[
 W_x\cdot\nabla f=h f
\]
for a polynomial $h$.  Every member of the fixed-$x$ family
through a real point of the regulus is then entirely contained
in it.  To verify this without an appeal to invariance of flows,
restrict to that line and write $g(t)=f(t,c(t))$.  The identity
becomes
\[
 (1+t^2)g'(t)=h(t,c(t))g(t).
\]
A nonzero polynomial $g$ cannot have a real root: at a root its
derivative has strictly smaller vanishing order, while $1+t^2$
does not vanish.  The chosen line has such a root, so $g$ is
identically zero.

Thus this fixed-$x$ family supplies a full real affine ruling of the
regulus.  Its lines are pairwise projectively skew, so they all
belong to the same ruling.  Conversely, coverage of every real
affine point implies that no real affine line of that ruling is
missing.  A ruling line contained in the plane at infinity is
irrelevant here.
Different fixed-$x$ families share no line, because the slope
and intercept recover both endpoints.  There are only two
rulings, so at most two values of $x$ contribute seven or more
lines.  They contribute at most $2v$ lines in all; the others
contribute at most $6u$.  This proves the claimed cap.
\end{proof}

\begin{theorem}[Planar distinct distances]\label{ci:planar-distances}
Every finite set $Q\subset\mathbb R^2$ of $n\ge2$ points determines
at least $c n/\log(2n)$ distinct positive distances.
\end{theorem}

\begin{proof}
Two distinct ordered matches determine a unique orientation-preserving
planar isometry whenever the matched segments have the same positive
length.  There are consequently finitely many direct isometries
with at least two matches from $Q$ to itself.  Rotate a second copy
of $Q$ by a generic rotation $R_0$ to form $Q'=R_0Q$.  We may
ensure that none of the direct isometries with two matches from
$Q$ to $Q'$ has linear part $-I$: each is $R_0g$ for one of the
finitely many just described, and each excludes only one rotation
$R_0$.  All these motions therefore lie in the Cayley chart used
in \eqref{ci:cayley-line}.  Translations lie in this chart as well.

Use the $n^2$ lines for the pairs $(x,y)\in Q\times Q'$.
Lemma~\ref{ci:cayley-caps} gives plane cap $n$ and regulus cap
$8n$.  Theorem~\ref{ci:two-rich} bounds the two-rich points by
$Cn^3$.  For $3\le k\le n$, Theorem~\ref{ci:rich-lines} bounds
the $k$-rich points by
\[
 C\left(\frac{n^3}{k^2}+\frac{n^3}{k^3}+\frac{n^2}{k}\right)
 \le C'\frac{n^3}{k^2}.
\]
Multiplicity equals the number of matches, by the bijective
Cayley representation, and it is at most $n$.

For each positive distance $\rho$ let $a_\rho$ be the number of
ordered pairs of distinct points of $Q$ at distance $\rho$.
The corresponding number for $Q'$ is also $a_\rho$.  Every
ordered pair of equally long nonzero segments, one from each
cloud, gives the unique direct isometry described above.  A motion
with $k$ matches contributes exactly $k(k-1)$ such pairs.
Splitting multiplicities into intervals $[2^j,2^{j+1})$, $j\ge1$,
and using the preceding bounds yields
\[
 \sum_\rho a_\rho^2
   =\sum_g k_g(k_g-1)
   \le C n^3\log(2n).
\]
Since $\sum_\rho a_\rho=n(n-1)$, Cauchy--Schwarz gives
\[
 |\Delta(Q)|\ge
 \frac{n^2(n-1)^2}{Cn^3\log(2n)}
 \ge c\frac n{\log(2n)}.
\]
\end{proof}

\begin{corollary}[Planar occupancy]\label{ci:planar-occupancy}
If a finite set in Euclidean space determines at most $M\ge1$ positive
distances, then any affine two-plane contains at most
$CM\log(2M)$ of its points.  In particular, along a sequence with
$B\longrightarrow\infty$ and $M=o(B^2)$, every two-plane has at
most $B^{2+o(1)}$ points, uniformly in the plane.
\end{corollary}

\begin{proof}
For an occupancy $n\ge2$, Theorem~\ref{ci:planar-distances} gives
$n\le CM\log(2n)$.  The elementary inequality
$\log(2n)\le C\sqrt n$ first yields $n\le C'M^2$; substitute
this back into the logarithm to get $n\le C''M\log(2M)$.
Occupancies zero and one are harmless.  If $M=o(B^2)$, eventually
$M\le B^2$, so the uniform bound is $O(B^2\log(2B))$, which is
$B^{2+o(1)}$.
\end{proof}
